% N4 --- Superposicao: 8 desafios, com circuitos quando pedagogicamente util
% Revisao visual preventiva: maior separacao entre ramos, fontes e rotulos.
\blocodesafio

\begin{questao}[4]{}
A rede linear abaixo é vista por um nó de saída. Demonstre, usando a forma matricial nodal, que a resposta total é a soma das respostas às fontes $E_1$, $I_2$ e $E_3$ separadamente.
\circuitoq{\begin{circuitikz}[american,scale=.68,transform shape,line width=.8pt]
\coordinate (g) at (5.8,0); \coordinate (a) at (5.8,4.5);
\draw (0,0) to[\fonteV,l^={$E_1$}] (0,4.5) to[R,l={$R_1$}] (a);
\draw (8.8,0) to[isource,l_={$I_2$}] (8.8,4.5)--(a);
\draw (12.0,0) to[\fonteV,l_={$E_3$}] (12.0,4.5) to[R,l_={$R_3$}] (a);
\draw (a) to[R,l_={$R_L$}] (g);
\draw (0,0)--(g)--(8.8,0)--(12.0,0);
\fill (a) circle (1.1pt); \node[Vcol,fill=white,inner sep=1.5pt] at ($(a)+(0,.68)$) {$V_o$};
\node[ground] at(g){};
\end{circuitikz}}
\resposta{$\mathbf v=\mathbf G^{-1}\mathbf b_1+\mathbf G^{-1}\mathbf b_2+\mathbf G^{-1}\mathbf b_3$}
\resolucao{Em uma rede linear, $\mathbf G$ não muda com o valor das fontes e $\mathbf b=\mathbf b_1+\mathbf b_2+\mathbf b_3$. Logo $\mathbf v=\mathbf G^{-1}\mathbf b=\sum_k\mathbf G^{-1}\mathbf b_k$. Cada termo é exatamente uma análise parcial.}
\end{questao}

\begin{questao}[4]{}
No resistor, as parcelas de corrente são $I_1$ e $I_2$. Expanda a potência e determine a condição para que a soma das potências parciais coincida com a potência real.
\circuitoq{\begin{circuitikz}[american,scale=.88,transform shape,line width=.8pt]
\coordinate (g) at (3.8,0); \coordinate (a) at (3.8,4.0);
\draw (0,0) to[isource,l^={$I_1$}] (0,4.0)--(a);
\draw (7.6,0) to[isource,l_={$I_2$}] (7.6,4.0)--(a);
\draw (a) to[R,l={$R$},i>^={$I$}] (g);
\draw (0,0)--(g)--(7.6,0);
\end{circuitikz}}
\resposta{$P=RI_1^2+RI_2^2+2RI_1I_2$; a soma parcial só coincide quando $I_1I_2=0$}
\resolucao{O termo cruzado $2RI_1I_2$ é a parcela ausente na soma ingênua. Ele pode ser positivo ou negativo conforme as contribuições tenham o mesmo sentido ou sentidos opostos.}
\end{questao}

\begin{questao}[4]{}
Projete um ensaio experimental de superposição no circuito, especificando as três configurações de medição e o modo correto de desativar cada fonte.
\circuitoq{\begin{circuitikz}[american,scale=.80,transform shape,line width=.8pt]
\coordinate (g) at (4.5,0); \coordinate (a) at (4.5,4.2);
\draw (0,0) to[\fonteV,l^={$E=\SI{15}{\volt}$}] (0,4.2) to[R,l={$\SI{5}{\ohm}$}] (a);
\draw (a) to[R,l_={$\SI{10}{\ohm}$}] (g)--(0,0);
\draw (8.8,0) to[isource,l_={$\SI{1}{\ampere}$}] (8.8,4.2)--(a); \draw (8.8,0)--(g);
\fill (a) circle (1.1pt); \node[Vcol,fill=white,inner sep=1.6pt] at ($(a)+(0,.68)$) {$V_A$};
\node[ground] at(g){};
\end{circuitikz}}
\resposta{Medir com ambas ativas, depois somente $E$, depois somente $I$; verificar $V_A\approx V_A^{(E)}+V_A^{(I)}$}
\resolucao{Ao medir apenas $E$, abre-se o ramo da fonte de corrente. Ao medir apenas $I$, substitui-se a fonte de tensão ideal por curto. Não se deve simplesmente retirar a fonte de tensão do circuito.}
\end{questao}

\begin{questao}[4]{}
Escolha $E_2$ para anular a corrente em $R$ sem alterar $E_1$.
\circuitoq{\begin{circuitikz}[american,scale=.90,transform shape,line width=.8pt]
\draw (0,0) to[\fonteV,l^={$E_1=\SI{30}{\volt}$}] (0,3.8)
      to[R,l={$R=\SI{10}{\ohm}$},i>^={$I$}] (5.8,3.8)
      to[\fonteV,l_={$E_2$},invert] (5.8,0)--(0,0);
\end{circuitikz}}
\resposta{$E_2=\SI{30}{\volt}$ com polaridade oposta à de $E_1$}
\resolucao{A parcela de $E_1$ é $+3$ A. Para corrente total nula, a segunda deve produzir $-3$ A; portanto $E_2=10(3)=30$ V em oposição.}
\end{questao}

\begin{questao}[4]{}
Com as fontes independentes zeradas, determine a resistência vista na porta aplicando uma fonte de teste, mantendo a fonte dependente ativa.
\circuitoq{\begin{circuitikz}[american,scale=.76,transform shape,line width=.8pt]
\coordinate (g) at (4.2,0); \coordinate (a) at (4.2,4.3);
\draw (a) to[R,l_={$\SI{1}{\kilo\ohm}$}] (g);
\draw (0,0) to[cisource,l^={$0{,}5\,\mathrm{mS}\,V_A$}] (0,4.3)--(a); \draw (0,0)--(g);
\draw (8.8,0) to[isource,l_={$I_t$}] (8.8,4.3)--(a); \draw (8.8,0)--(g);
\fill (a) circle (1.1pt); \node[Vcol,fill=white,inner sep=1.6pt] at ($(a)+(0,.68)$) {$V_A$};
\node[ground] at(g){};
\end{circuitikz}}
\resposta{$R_{eq}=\SI{2}{\kilo\ohm}$}
\resolucao{$I_t=V_A/1000-0{,}0005V_A=0{,}0005V_A$. Assim $V_A/I_t=2000\,\Omega$. A dependente não pode ser desativada.}
\end{questao}

\begin{questao}[4]{}
Uma rede tem quatro fontes, mas deseja-se apenas saber qual delas domina $V_o$. Proponha uma estratégia de superposição que minimize o trabalho e produza uma medida quantitativa de dominância.
\circuitoq{\begin{circuitikz}[american,scale=.58,transform shape,line width=.8pt]
\coordinate (g) at (7.0,0); \coordinate (a) at (7.0,5.0);
\draw (0,0) to[\fonteV,l^={$E_1$}] (0,5.0) to[R,l={$R_1$}] (a);
\draw (2.5,0) to[isource,l^={$I_2$}] (2.5,5.0)--(a);
\draw (11.5,0) to[isource,l_={$I_4$}] (11.5,5.0)--(a);
\draw (14.0,0) to[\fonteV,l_={$E_3$}] (14.0,5.0) to[R,l_={$R_3$}] (a);
\draw (a) to[R,l_={$R_L$}] (g);
\draw (0,0)--(2.5,0)--(g)--(11.5,0)--(14.0,0);
\fill (a) circle (1.1pt); \node[Vcol,fill=white,inner sep=1.5pt] at ($(a)+(0,.72)$) {$V_o$};
\node[ground] at(g){};
\end{circuitikz}}
\resposta{Calcular as quatro parcelas $V_{o,k}$ e comparar $|V_{o,k}|/\sum_j|V_{o,j}|$}
\resolucao{A fatoração da matriz nodal pode ser reutilizada nas quatro excitações. O índice normalizado pelos módulos evita que cancelamentos algébricos escondam uma fonte individualmente forte.}
\end{questao}

\begin{questao}[4]{}
Formule um algoritmo para calcular simultaneamente $n$ saídas de uma rede linear com $m$ fontes, evitando refatorar a matriz a cada parcela.
\resposta{Fatorar $\mathbf G$ uma vez; resolver $m$ vetores de excitação; armazenar as respostas em uma matriz de transferência $\mathbf H$}
\resolucao{Monte $\mathbf G=\mathbf L\mathbf U$ uma única vez. Para cada fonte $k$, resolva $\mathbf G\mathbf v_k=\mathbf b_k$ por substituição. As $n$ grandezas selecionadas formam a coluna $k$ de $\mathbf H$, permitindo depois $\mathbf y=\mathbf H\mathbf s$.}
\end{questao}

\begin{questao}[4]{}
No circuito, indique até que faixa de operação a superposição é válida e qual evento físico faz o modelo deixar de ser linear.
\circuitoq{\begin{circuitikz}[american,scale=.74,transform shape,line width=.8pt]
\coordinate (g) at (4.4,0); \coordinate (a) at (4.4,4.2);
\draw (0,0) to[\fonteV,l^={$E_1$}] (0,4.2) to[R,l={$R_1$}] (a);
\draw (8.8,0) to[\fonteV,l_={$E_2$}] (8.8,4.2) to[R,l_={$R_2$}] (a);
\draw (a) to[R,l_={$R_L$}] (g); \draw (0,0)--(g)--(8.8,0);
\fill (a) circle (1.1pt); \node[Vcol,fill=white,inner sep=1.5pt] at ($(a)+(0,.68)$) {$V_o$};
\node[draw=cinza,rounded corners,align=center,font=\footnotesize,inner sep=3pt]
      at (11.4,2.1) {fonte real com\\limitação de corrente};
\node[ground] at(g){};
\end{circuitikz}}
\resposta{Válida enquanto resistores e fontes operarem em região linear; falha quando uma fonte entra em limitação/saturação ou um elemento muda de estado}
\resolucao{A superposição não é uma regra de desenho, mas consequência da linearidade do modelo. Se a fonte satura, a relação entre excitação e resposta deixa de ser proporcional e as parcelas já não podem ser somadas.}
\end{questao}